\section{Retrieval + Reasoning：先找到相关结构，再展开计算}

有些失败不是模型不会做推理，而是输入缺少事实、公式、先例或更高层抽象。Retrieval 与 reasoning 解决不同环节：前者改变模型可见的信息，后者在给定信息上执行分解与计算。课堂拒绝“二选一”争论，直接主张 retrieval + reasoning，因为产品目标是整体正确率而不是维护方法阵营。

\subsection{为什么检索与推理不是替代关系}

slide 43 的问题只有一行，却总结了系统分层：检索器决定把什么证据放入上下文，推理器决定如何组合证据并得到结论。检索命中但推理错误，会出现引用正确却结论错误；推理正确但检索缺失，会在错误前提上严谨演算。两者需要分别评测 recall、provenance、entailment 与最终任务成功。

\lecturefigure{slide-43.jpg}{Retrieval or reasoning？课堂答案是 Retrieval + Reasoning}{43}

设检索器从语料库 $\mathcal C$ 返回 $z\sim q_\phi(z\mid x)$，生成器再输出 $y\sim p_\theta(y\mid x,z)$。理想答案概率为
\[
p(y\mid x)=\sum_{z\in\mathcal C}q_\phi(z\mid x)p_\theta(y\mid x,z).
\]
其中 $\phi$ 是检索参数，$\theta$ 是生成参数，$z$ 是证据。实际系统只取 top-$k$ 文档近似求和；因此 $k$、重排器、上下文压缩和引用保留都会改变结果。检索不是把网页文本塞进 prompt 就结束，而是一个有召回率、噪声和时效性的概率组件。

\begin{knowledgebox}{首次使用：Retrieval-Augmented Reasoning}
它指先从外部文档、记忆、工具或案例库取得相关信息，再让模型基于证据进行中间步骤与答案生成。必须保留来源和时间戳，区分“模型记得”“检索到了”“由证据推导”，否则错误无法定位。
\end{knowledgebox}

\subsection{Analogical Reasoning：检索相关问题而非相同答案}

slide 44 给出一个旋转正方形：四个顶点为 $(-2,2)$、$(2,-2)$、$(-2,-6)$、$(-6,-2)$。直接问面积时模型可能失败；提示 “recall a related problem” 后，它先召回两点距离公式。相邻顶点距离为
\[
s=\sqrt{(2-(-2))^2+(-2-2)^2}=\sqrt{32},
\qquad A=s^2=32.
\]
其中 $s$ 是正方形边长，$A$ 是面积。被检索的不是同一道题，而是可迁移的子结构“两点距离”；这说明 retrieval 可以返回算法原语、证明模板或相似失败案例，而不只返回事实段落。

\lecturefigure{slide-44.jpg}{Analogical reasoning：先召回相关距离问题，再求旋转正方形面积}{44}

\teachervoice{00:51:29--00:55:00，老师说自己在工业界更关心性能，不想争论 retrieval 还是 reasoning。课堂提示：related problem 不等于 duplicate problem；有用检索往往提供可迁移原理，而不是直接泄漏答案。}

\begin{warningbox}{类比检索的负迁移}
表面相似案例可能使用不同约束、单位或因果机制。系统应让模型明确写出“哪些结构相同、哪些条件不同”，并在应用类比前验证前提；否则相关性检索会把最像的错误模板带进推理。
\end{warningbox}

\subsection{Step-Back：先抽象，再回到细节}

slide 45 的 step-back prompting 不是简单“想更久”，而是把原问题提升到更一般的概念或原理：物理题先找守恒定律，时间约束问答先整理完整经历，再回到具体区间。这个操作把检索 query 从表面措辞改写为结构性问题，常能找到更稳定的知识。

\lecturefigure{slide-45.jpg}{Step-Back Prompting：通过抽象召回原则，再解决原问题}{45}

可把流程写成三步：
\begin{enumerate}
  \item 生成抽象问题 $x' = g(x)$，例如“这类系统遵循什么守恒关系”；
  \item 检索或生成原则 $z=\operatorname{retrieve}(x')$；
  \item 在原约束下求解 $y=\operatorname{reason}(x,z)$。
\end{enumerate}
关键验收是 $z$ 是否真的适用于 $x$。若抽象过宽，检索会返回常识废话；若过窄，就退化成重述原题。可以记录抽象 query、证据命中和最终引用，让失败分析知道问题出在 query rewrite、retrieval 还是 reasoning。

\subsection{Deep Research 是多轮检索—推理循环}

slide 46 用 Deep Research 产品入口说明同一思想如何扩展到长任务：系统不只检索一次，而是形成子问题、搜索、阅读、更新假设、继续检索，最后组织带来源的报告。课堂画面同时出现 Gemini Deep Research 标题与搜索/Deep research 产品入口，教学重点是产品类别，不应把 UI 当作 2026 年功能保证。

\lecturefigure{slide-46.jpg}{Deep Research：把检索、分解与综合组织成多轮任务}{46}

\begin{lstlisting}[language=Python,caption={带 provenance 与停止条件的 retrieval-reasoning 循环}]
def deep_research(question, search, reasoner, budget):
    state = {"open": [question], "evidence": [], "claims": []}
    while state["open"] and budget.remaining():
        subquestion = state["open"].pop(0)
        documents = search(subquestion, top_k=8)
        update = reasoner.synthesize(subquestion, documents)
        state["evidence"].extend(update.citations)
        state["claims"].extend(update.claims)
        state["open"].extend(update.followups)
        budget.charge(documents, update)
    return reasoner.final_report(state, require_citations=True)
\end{lstlisting}

这个循环必须有停止条件和证据去重。没有预算上限，系统可能不断提出新问题；没有 provenance，最终文本无法核查；没有 claim--evidence 对齐，引用只是装饰。对于时效性事实还要保存抓取日期，因为同一 URL 的内容可能变化。

\teachervoice{00:55:00--00:56:31，老师把 analogical reasoning、step-back 和 deep research 连接起来：先找到相关知识或更高层结构，再做推理。讲义提醒：真实 deep research 还必须处理来源质量、冲突、时效和终止，不是多搜几次就自动可靠。}

\subsection{本章小结}

Retrieval 为 reasoning 提供外部事实、案例与抽象原理；reasoning 把这些证据转成答案。Analogical prompting、step-back 和 deep research 可以看作从单次类比到多轮研究的连续谱。系统质量必须分解为检索召回、证据可靠性、推理正确性与引用忠实性。

